% Avhengigheits-sloyfa (Eyal sin «Hook»): Utloysar -> Handling -> Variabel
% belonning -> Investering -> Utloysar. Under: Fogg B = MAT.
% Reint TikZ, ingen ekstra bibliotek. Input-bart fragment.
\begin{figure}[t]
\centering
\begin{tikzpicture}[
  font=\footnotesize,
  boks/.style={draw, rounded corners=2pt, align=center, inner sep=4pt, text width=30mm, fill=black!4},
  pil/.style={-{Latex[length=2mm]}, line width=0.7pt},
]
% Fire nodar i ein loop (klokka rundt), radius ca 2.6
\node[boks] (ut)  at (0,2.6)   {\textsc{Utløysar}\\\scriptsize ekstern eller intern};
\node[boks] (ha)  at (3.0,0)   {\textsc{Handling}\\\scriptsize det enklaste steget};
\node[boks] (be)  at (0,-2.6)  {\textsc{Variabel belønning}\\\scriptsize uforutsigbar gevinst};
\node[boks] (in)  at (-3.0,0)  {\textsc{Investering}\\\scriptsize brukaren legg noko inn};

% Piler langs sloyfa (med klokka)
\draw[pil] (ut.east)  to[bend left=18] (ha.north);
\draw[pil] (ha.south) to[bend left=18] (be.east);
\draw[pil] (be.west)  to[bend left=18] (in.south);
\draw[pil] (in.north) to[bend left=18] (ut.west);

% Senter-merknad
\node[align=center, font=\scriptsize, text width=26mm] at (0,0)
  {sløyfa lukkar\\seg om seg sjølv};

% Fogg-formelen under
\node[draw, rounded corners=2pt, anchor=north, text width=72mm, align=center, inner sep=5pt]
  at (0,-4.0)
  {\textsc{Fogg:}\quad $B = MAT$\\[2pt]\scriptsize åtferd $=$ motivasjon $\times$ evne $\times$ utløysar};
\end{tikzpicture}
\caption{\textsc{Avhengigheits-sløyfa.} Eyals «hook»-modell formaliserer korleis ein produkt-sløyfe vert vane: ein utløysar driv ei handling, ei variabel belønning held ho ved like, og ei investering gjer neste runde sannsynlegare. Foggs åtferdsmodell ($B = MAT$) gjev oppskrifta på kvart steg: senk terskelen for handling og plassér utløysaren der motivasjonen alt er. Dette er ikkje ein teori om god form, men ei teknikk for å gjere bruk til tvang.}
\end{figure}
